\documentclass[10pt,a4paper]{article}
\usepackage[margin=26mm]{geometry}
\usepackage{amsmath,amssymb,lmodern}
\setlength{\emergencystretch}{2em}
\title{Firm nonexpansiveness and fixed points of resolvents\\
\large Proof of Conjecture 00000008844}
\author{ziangni-sys}
\date{4 October 2026}
\begin{document}
\maketitle

\section*{Definitions and domain}
Let $H$ be a real Hilbert space and let $A:H\rightrightarrows H$ be
monotone: for every $a\in A(u)$ and $b\in A(v)$,
\[
\langle u-v,a-b\rangle\geq0.
\]
For the standard positive parameter $\lambda>0$, define the resolvent
$J_\lambda=(I+\lambda A)^{-1}$ by
\[
u\in J_\lambda(x)
\quad\Longleftrightarrow\quad
\exists a\in A(u):\ x=u+\lambda a.
\]
Its natural domain is $D_\lambda=\operatorname{ran}(I+\lambda A)$.
We will prove that its values there are singletons, so it is an actual
map $J_\lambda:D_\lambda\to H$.
The source conjecture's two properties hold on this natural domain.
No claim of totality on $H$ is made: monotonicity alone does not give
that additional property. For example, the graph containing only $(0,0)$
is monotone and its resolvent domain is $\{0\}$.

\section*{Firm nonexpansiveness}
Take $x,y\in D_\lambda$, $u\in J_\lambda(x)$, and $v\in J_\lambda(y)$.
Choose $a\in A(u)$ and $b\in A(v)$ with
$x=u+\lambda a$, $y=v+\lambda b$. Then
\begin{align*}
\langle u-v,x-y\rangle
&=\langle u-v,(u-v)+\lambda(a-b)\rangle\\
&=\|u-v\|^2+\lambda\langle u-v,a-b\rangle
\geq\|u-v\|^2.
\end{align*}
If $x=y$, the inequality forces $\|u-v\|^2=0$, hence $u=v$.
Existence on $D_\lambda$ follows from its definition, and thus the
resolvent is single-valued there. Substituting its unique values yields
\[
\|J_\lambda(x)-J_\lambda(y)\|^2
\leq\langle J_\lambda(x)-J_\lambda(y),x-y\rangle,
\]
which is the defining firm nonexpansiveness inequality.

\section*{Fixed points and zeros}
For any $x\in H$, the diagonal graph relation is equivalent to a zero:
\[
x\in J_\lambda(x)
\quad\Longleftrightarrow\quad
\exists a\in A(x):\ x=x+\lambda a
\quad\Longleftrightarrow\quad
0\in A(x).
\]
The second equivalence uses $\lambda>0$, hence $\lambda a=0$ implies $a=0$.
A zero therefore automatically belongs to $D_\lambda$. By single-valuedness,
the displayed equivalence proves the exact set identity
\[
\{x\in D_\lambda:J_\lambda(x)=x\}
=\{x\in H:0\in A(x)\}.
\]

\section*{Formal verification}
Lean defines the set-valued graph, monotonicity, inverse graph, natural
domain, and a choice-defined resolvent map. It proves the map satisfies
the inverse graph and that this graph is exactly its value relation.
The theorems \texttt{resolvent\_firm} and
\texttt{fixed\_set\_eq\_zero\_set} establish both original clauses.
All results are for arbitrary real inner-product spaces, so they include
Hilbert spaces without needing completeness or maximality.
\end{document}
